\section{Introduction}

Internal ribosome entry sites (IRESes) enable cap-independent translation and
are useful regulatory components in RNA engineering. Their function is not
defined by a short motif alone: sequence context, alternative RNA structures,
host factors, cargo, and assay architecture can alter the measured output.
Consequently, a design procedure that only maximizes an IRES-recognition score
can produce candidates whose apparent improvement is disconnected from the
structural state of an experimentally supported parent.

Recent deep-learning work established IRES identification and evolutionary
sequence optimization baselines using 46,774 labelled sequences
\cite{chu2026iresai}. The dominant source, however, is a 174-nt
DNA/lentiviral reporter library \cite{weingarten2016ires}; it is not identical
to full-length or direct-RNA settings. We reproduced the released RNA-FM and
UTR-LM classifier summaries on their native repeated holdouts
(AUROC/AUPR/F1: $0.778/0.614/0.510$ and $0.765/0.581/0.493$,
respectively), but an all-checkpoint RNA-FM ensemble reached only AUROC 0.600
on a zero-shot direct-RNA evaluation. This motivates treating a historical
recognition model as a proposal signal rather than a complete activity oracle.

We present \textsc{StructIRES}, a two-stage framework linking structure-aware
IRES recognition with constrained full-length IRES optimization. Its
\textsc{StructIRES-Classifier} is a validation-clean RNA-FM re-training with
an MFE-contact encoder and gated residual; \textsc{StructIRES-Rank} is the
full-length constraint-guided selection component. The framework makes four
contributions:
\begin{enumerate}
  \item It tests whether an RNA foundation-model sequence prior is improved by
  contextual MFE base-pair features under matched validation-clean native
  folds, while reporting a separate identity-aware robustness audit.
  \item It combines sequence plausibility with parent-relative energy,
  thermodynamic ensemble, and ensemble-anchor preservation during design.
  \item It evaluates design alternatives on identical frozen candidate pools and
  selection budgets, isolating the value of each constraint.
  \item It uses direct-RNA evidence only as an explicitly assay-qualified
  evaluator or a parent-specific mutational guardrail, never as an inferred
  activity label for newly generated candidates.
\end{enumerate}

\begin{figure*}[t]
  \centering
  \includegraphics[width=0.96\textwidth]{figures/fig1_overview.pdf}
  \caption{Overview of \textsc{StructIRES}. \textbf{A}, the recognition
  component augments an author-style RNA-FM classifier with contextual
  residue-pair features over label-free ViennaRNA MFE contacts, fused through
  a zero-initialized gated residual. \textbf{B}, the completed design study
  ranks a shared, length-preserving mutation pool using a frozen released
  IRES-LM proposal score plus parent-relative energy, ensemble-profile and
  anchor-preservation objectives. \textbf{C}, direct-RNA resources are kept
  outside selection: they serve only as a held-out computational proxy or an
  IAPV-specific MPRA edit-risk guardrail.}
  \label{fig:overview}
\end{figure*}
